\clearpage\begingroup
\hypertarget{native-307-lower}{}\HistoricalHead{302}{L}{PRAKTISCHE UND SPHÄRISCHE ASTRONOMIE.}
\begin{center}[II.]\end{center}
\noindent\hspace*{\parindent}Unterthänigstes P. M.\\
Die Bestellung der astronomischen\\
Instrumente für die neue Sternwarte\\
in Göttingen betreffend.

Da jetzt die Zeit eingetreten ist, wo wiederum an dem Baue der hiesigen neuen Sternwarte fortgefahren wird, und die gänzliche Vollendung dieses Baues im gegenwärtigen Jahre zu hoffen steht, so scheint es nunmehro Zeit zu sein, Maassregeln wegen Ausrüstung der Sternwarte mit denjenigen Instrumenten zu treffen, durch welche sie erst zu dem ihr gebührenden Range erhoben werden wird. Ich will mich hier in kein weitläuftiges Detail wegen dieser
\clearpage\SourceLeaf{308}\HistoricalHead{303}{R}{BERICHTE UND ANTRÄGE ZUR ANSCHAFFUNG VON INSTRUMENTEN.}
\noindent Instrumente einlassen, da ich mich in dieser Hinsicht auf ein früheres unterthäniges P. M. beziehen kann, und also nur kurz bemerken, dass, nachdem der Ankauf des \textsc{Repsold}schen Instruments genehmigt worden (welches im Laufe dieses Sommers abgeliefert werden soll) und nachdem die reichhaltige Sammlung von Sehwerkzeugen aus Lilienthal grösstentheils schon hier angelangt ist, unsre \textls[120]{Hauptbedürfnisse} in denjenigen fixen Instrumenten bestehen werden, von welchen ich in vorgedachtem P. M. umständlicher gesprochen habe, nemlich
\begin{center}
\(\left.\begin{tabular}{@{}rl@{}}1)&einem Mittagsfernrohre\\2)&einem Meridiankreis\end{tabular}\right\}\) beide im westlichen Observationszimmer,\par\smallskip
\(\left.\begin{tabular}{@{}rp{.68\textwidth}@{}}3)&einem Aequatoreal, oder\\&einem grossen parallatisch aufzustellenden Refraktor\end{tabular}\right\}\) in der Rotunde.
\end{center}

Solange die Sternwarte diese Instrumente, und besonders die beiden ersten, noch nicht besitzt, wird sie ihre Wirksamkeit auf eine ganz würdige und vollständige Art noch nicht beginnen können, und es scheint daher um so nothwendiger, diese Instrumente jetzt sobald als möglich zu bestellen, da ohnehin von der Bestellung bis zur Vollendung, Ablieferung und Aufstellung derselben immer noch eine geraume Zeit verstreichen wird.

Die erste Frage ist nun, von welchem Künstler diese Instrumente in grösster Vollkommenheit erwartet werden können. Es scheint allgemein anerkannt zu sein, dass die Arbeiten des grossen Künstlers \textsc{Reichenbach} in München von den Arbeiten keines andern jetzt lebenden Künstlers übertroffen werden, und ich glaube daher, dass unsre Sternwarte sich Glück wünschen kann, wenn sie von jenem Künstler ihre Ausrüstung erhält. Allein es ist nicht hinreichend, einem Künstler aufzugeben, ein Instrument von dieser oder jener Benennung in vorgeschriebener Grösse anzufertigen. Bei Nr. 1, dem Mittagsfernrohr, welches seiner Natur nach das am wenigsten zusammengesetzte ist, liesse sich dies allenfalls thun. Allein die Instrumente, welche den Namen Meridiankreis führen, sind, obgleich sie einen und denselben letzten Zweck haben, in den letzten 25 Jahren von so mannichfaltiger Einrichtung gebaut, und jede Einrichtung hat so manche eigenthümliche Vortheile, dass es unter den Astronomen noch unentschieden ist, welcher Gattung von Kreisen der Preis der grössten Vortrefflichkeit beizulegen sei. Die Meridiankreise, mit welchen \textsc{Piazzi} in Palermo, \textsc{Pond} in Greenwich, \textsc{Oriani} und \textsc{Bouvard} in Mayland und
\clearpage\SourceLeaf{309}\HistoricalHead{304}{L}{PRAKTISCHE UND SPHÄRISCHE ASTRONOMIE.}
\noindent Paris, \textsc{Bessel} in Königsberg observiren (resp. von \textsc{Ramsden}; \textsc{Troughton}; \textsc{Reichenbach}; \textsc{Carry}) sind gewissermaassen ganz verschiedene Instrumente, und eben jetzt ist \textsc{Reichenbach} wiederum mit einer ganz neuen Art beschäftigt, nach welcher er Kreise, für Ofen, Turin, Warschau und München bestimmt, construiren will. — Eine ähnliche, wenn auch nicht ganz so grosse, Mannichfaltigkeit findet auch bei No. 3, dem Aequatoreal, statt.

Unter dieser Umständen scheint in der That nichts wünschenswerther und nothwendiger, als eine mündliche Rücksprache des Astronomen mit dem Künstler, wo beide ihre Ideen gegenseitig austauschen, Vortheile und Nachtheile gegen einander abwägen, und wo die Bedürfnisse der Wissenschaft, die Bedenklichkeiten, welchen diese oder jene Einrichtungen ausgesetzt sein möchten, die Hülfsmittel, welche das Genie des Künstlers findet, die Schwierigkeiten zu überwinden, gehörig erwogen und ausgeglichen werden können, nicht zu gedenken, dass die persönliche Communication des Astronomen mit dem Künster auch im Allgemeinen für die künftige Aufstellung, Behandlung und Schonung der Instrumente, wenn er deren Bau in der Werkstatt des Künstlers selbst gesehen hat, von dem allergrössten Nutzen sein muss.

Auch in anderen ähnlichen Fällen hat man eine solche Maassregel für dienlich gehalten. So wurde vor Erbauung der Palermer Sternwarte \textsc{Piazzi} von der Neapolitanischen Regierung nach England geschickt, um bei \textsc{Ramsden} die Instrumente zu bestellen, und hielt sich dort selbst mehrere Jahre auf, um die Vollendung zu betreiben (welches letztere glücklicherweise bei \textsc{Reichenbach} nicht nöthig sein wird, wenn er einmal die Verfertigung übernommen haben wird); gleicherweise wurden \textsc{Pasquich} von Ofen, \textsc{Tralles} von Berlin nach München geschickt, und noch mehrere ähnliche Beispiele liessen sich anführen, wenn es nöthig wäre.

Sollte königl. Curatorium diesen Ansichten Beifall schenken, und für gut finden, mich mit dem Auftrage der Reise nach München und der persönlichen Communication mit Herrn \textsc{Reichenbach} zu beehren, so scheint mir die Zeit der bevorstehenden Osterferien dazu am gelegensten, wo, eingezogener Erkundigung zu Folge, \textsc{Reichenbach} in München anwesend sein wird. Ich würde dann suchen, diesen grossen Künstler zur Uebernahme der Verfertigung unsrer Bedürfnisse zu disponiren, alles Nöthige mit ihm verabreden, und nach meiner Zurückkunft dem königl. Curatorium von den Resultaten der Verabredung,
\clearpage\SourceLeaf{310}\HistoricalHead{305}{R}{BERICHTE UND ANTRÄGE ZUR ANSCHAFFUNG VON INSTRUMENTEN.}
\noindent den Bedingungen, unter welchen der Künstler die Instrumente liefern will u. s. w. Bericht abstatten, damit sodann die definitive Bestellung vom k. Curatorium verfügt werden könne, und die hiesige neue Sternwarte die gewisse Aussicht erhalte, sobald als möglich mit den andern vom ersten Range ehrenvoll wetteifern zu können.
\begin{flushright}Unterthänigst\end{flushright}
Göttingen den 24. März 1816.
\begin{flushright}\textsc{Carl Friedrich Gauss}.\end{flushright}
\clearpage\endgroup
